\chapter{Avance del objetivo específico 3}

\section{Protocolo de evaluación}

La evaluación compara el DINOv2 público con los brazos adaptados A* y B*. F1 macro es la métrica primaria; exactitud balanceada, AUC, ECE y sensibilidad de melanoma complementan la lectura. Las comparaciones confirmatorias deben fijar de antemano el conjunto de datos, arquitectura, preprocesamiento, particiones, semillas y regla de decisión.

\begin{figure}[htbp]
\centering
\begin{tikzpicture}[
    box/.style={draw, rounded corners, align=center, text width=8.5cm, minimum height=0.85cm, inner sep=5pt},
    phase/.style={box, fill=gray!10},
    external/.style={box, dashed, fill=orange!8},
    flow/.style={-{Stealth[length=2mm]}, thick},
    node distance=0.35cm
]
\node[phase] (manifest) {Congelar manifiestos\\IDs, hashes y exclusiones};
\node[phase, below=of manifest] (split) {Separar por lesión en HAM10000, manteniendo \texttt{lesion\_id} dentro de una sola partición};
\node[phase, below=of split] (fit) {Adaptar el codificador SSL y ajustar el clasificador lineal};
\node[phase, below=of fit] (select) {Seleccionar variantes solo con el conjunto interno de selección};
\node[phase, below=of select] (final) {Evaluar según el protocolo registrado en la partición de validación};
\draw[flow] (manifest) -- (split);
\draw[flow] (split) -- (fit);
\draw[flow] (fit) -- (select);
\draw[flow] (select) -- (final);
\node[external, below=of final] (pad) {PAD-UFES-20: evaluación externa de dominio, sin ajuste del modelo};
\node[phase, below=of pad] (metrics) {Reportar métricas, incertidumbre y límites por dominio};
\draw[flow, dashed] (final) -- (pad);
\draw[flow] (pad) -- (metrics);
\end{tikzpicture}
\caption{Flujo de evaluación y separación entre selección interna y evaluación externa.}
\label{fig:protocolo-evaluacion}
\end{figure}
\noindent\textit{Nota.} El diagrama representa el protocolo confirmatorio propuesto. Las lecturas v7 y v7.1 incluyeron evaluaciones previas de validación; v7.1 es una lectura secundaria adicional registrada después de seleccionar variantes en el conjunto interno de selección.

\section{Resultados internos y reproducibilidad estadística}

Las tres semillas v7 muestran que B* cumple la predicción registrada P3-IC en el contraste con DINOv2 público. A* cumple la regla primaria P1 en las tres semillas, pero no cumple el criterio secundario estricto P1-NI de no inferioridad en la semilla 42; por tanto, P1-NI no se declara robusta a la semilla. La diferencia B*--A* no ofrece evidencia consistente de que el término de tono mejore la clasificación.

Los resultados publicados se reportan según su protocolo y configuración, pero sus intervalos en HAM10000 remuestrean imágenes. Dado que \texttt{lesion\_id} agrupa imágenes correlacionadas de una misma lesión, los intervalos actuales no representan incertidumbre agrupada por lesión. El documento debe presentar esa limitación y reservar los intervalos agrupados como reanálisis pendiente, sin reemplazar retrospectivamente los resultados registrados.

\section{Evaluación externa de dominio}

PAD-UFES-20 contiene fotografías clínicas de teléfono móvil, no dermatoscopía. En su lectura primaria (BCC, melanoma y nevo; 1.141 imágenes de 720 pacientes), el DINOv2 público obtuvo F1 macro 0,424. Los brazos sin emparejamiento por tono obtuvieron 0,339 y 0,316, y los brazos con emparejamiento por tono 0,317 y 0,293 para las semillas 42 y 43, respectivamente; los cuatro resultados fueron inferiores al control público en los contrastes pareados de F1 macro. Esta evaluación muestra que la ventaja observada en HAM10000 no se transfiere automáticamente a imágenes clínicas de otro dominio. El remuestreo de esta evaluación se agrupó por paciente, según el protocolo registrado \parencite{pacheco2020padufes}.

Los metadatos de PAD-UFES-20 no permiten concluir sobre fototipos V--VI: esos grupos tienen muy pocas observaciones. Además, el ITA del proyecto no es un estimador Fitzpatrick validado. Por estas razones, los resultados no se presentan como una auditoría concluyente de equidad por fototipo.

\section{Estado del objetivo y trabajo pendiente}

Existe evidencia de evaluación interna y una prueba externa de dominio, pero el objetivo 3 no se considera cerrado. Restan: (i) comparar SSL y la línea base supervisada con el mismo conjunto de datos, arquitectura y partición por lesión; (ii) recalcular la incertidumbre de HAM mediante remuestreo bootstrap agrupado por lesión; (iii) verificar la procedencia exacta de los conjuntos de preentrenamiento; y (iv) mantener la evaluación externa libre de ajuste. Ninguna de estas tareas autoriza por sí sola nuevas corridas; cualquier corrida se ejecuta en UNAB con presupuesto previamente acordado.
